\begin{lemma}[Expectation-to-PAC conversion at the ALMM constants]\label{lem:step-003-markov-bridge}
Under the primitive one-block definitions, fix integers \(N\ge2\) and
\(M\ge1\), and let

\[
B:([N]\times\{0,1\})^M\to\{0,1\}^{[N]}
\]

be any randomized unrestricted map. If, for every \(t\in[N+1]\) and every
probability law \(Q\) on \([N]\),

\[
\mathbb E_{\substack{S\sim(Q^{\tau_t})^M\\g\sim B(S)}}
  R_Q(g,\tau_t)
\le 2^{-8},
\tag{1}
\]

then, for every such \(t,Q\),

\[
\Pr_{\substack{S\sim(Q^{\tau_t})^M\\g\sim B(S)}}
  \left[R_Q(g,\tau_t)>\frac1{16}\right]
\le\frac1{16}.
\tag{2}
\]

Thus \(B\) satisfies (W-PAC).
\end{lemma}

\begin{proof}
Fix an arbitrary \(t\in[N+1]\) and arbitrary law \(Q\) on \([N]\), while
keeping the same fixed sample size \(M\). Under the joint randomness of
\(S\sim(Q^{\tau_t})^M\) and \(g\sim B(S)\), define the proof-local random
variable

\[
Z_{t,Q}:=R_Q(g,\tau_t).
\]

For every realized output \(g\), population 0-1 risk is a probability, so

\[
0\le Z_{t,Q}\le1.
\tag{3}
\]

Apply Markov's inequality at the positive threshold \(1/16\). The strict
failure event in (W-PAC) is contained in the weak-threshold event, hence

\[
\begin{aligned}
\Pr\!\left(Z_{t,Q}>\frac1{16}\right)
&\le \Pr\!\left(Z_{t,Q}\ge\frac1{16}\right)\\
&\le 16\,\mathbb E Z_{t,Q}\\
&\le 16\cdot2^{-8}\\
&=2^{-4}
=\frac1{16}.
\end{aligned}
\tag{4}
\]

The final inequality remains valid when the expectation in (1) equals
\(2^{-8}\); no strict inequality has been inserted. Since \((t,Q)\) was
arbitrary, (4) holds simultaneously as a universally quantified statement
for every fixed \((t,Q)\), which is exactly (W-PAC). Randomization and
improperness of \(g\) do not affect (3). \end{proof}

\begin{proposition}[Expected-loss hardness below the one-block threshold bound]\label{prop:step-003-expected-hardness}
Let \(b_*,d_*>0\) and \(N_*\ge2\) be the constants supplied by the Proposition~\ref{prop:step-002-wrapper}. Under that proposition and
Lemma~\ref{lem:step-003-markov-bridge}, for every pair of integers
\(N\ge N_*\), \(M\ge8\) satisfying

\[
M<b_*\log_2^*N,
\]

every randomized unrestricted map

\[
B:([N]\times\{0,1\})^M\to\{0,1\}^{[N]}
\]

that is
\(\left(0.1,d_*/(M^2\log M)\right)\)-DP under one-row replacement admits
some \(t\in[N+1]\) and some probability law \(Q\) on \([N]\) for which

\[
\mathbb E_{\substack{S\sim(Q^{\tau_t})^M\\g\sim B(S)}}
  R_Q(g,\tau_t)
>2^{-8}.
\tag{5}
\]
\end{proposition}

\begin{proof}
Fix \(N,M\) and an arbitrary randomized unrestricted \(B\) satisfying all
the displayed hypotheses. Suppose (5) were false. The exact logical
negation of the existential strict inequality is

\[
\text{for every }t\in[N+1]\text{ and every law }Q\text{ on }[N],
\qquad
\mathbb E_{\substack{S\sim(Q^{\tau_t})^M\\g\sim B(S)}}
  R_Q(g,\tau_t)
\le2^{-8}.
\tag{6}
\]

Lemma~\ref{lem:step-003-markov-bridge} applies to the same map \(B\), the
same exact sample size \(M\), and every fixed \((t,Q)\). It converts (6) into
(W-PAC), including the allowed equality
\(\Pr(R_Q(g,\tau_t)>1/16)=1/16\).

All remaining hypotheses of Proposition~\ref{prop:step-002-wrapper} now hold:
\(N\ge N_*\), integer \(M\ge8\), the map is arbitrary and improper, and its
privacy parameters equal the permitted endpoints

\[
\varepsilon=0.1,
\qquad
\delta=\frac{d_*}{M^2\log M}.
\]

The denominator is positive even at \(M=8\), since \(\log8>0\). The wrapper therefore yields

\[
M\ge b_*\log_2^*N,
\]

contradicting the strict hypothesis \(M<b_*\log_2^*N\). Hence (6) is false,
so a fixed pair \((t,Q)\) satisfying (5) exists. The pair may depend on the
fixed learner \(B\), because only worst-instance hardness is asserted here; no
learner-independent prior or minimax conclusion is asserted. \end{proof}
